% method.tex — High-level ICLR method; implementation details belong in Appendix B.2.

\section{Method}
\label{sec:method}

As illustrated in \cref{fig:method-overview}, \ours{} learns a joint
representation of fitted coordinates and experimental electron density. We
voxelize each crystallographic protein--ligand instance into aligned atomic and
density channels, and pre-train a Channel-ViT by reconstructing masked spatial
regions. The frozen encoder is then transferred through a pooled probe for
affinity regression or a spatial conditioning module for molecular generation.

\input{NeuripsW2026-AI4DD/resource/figure/method-overview}

\subsection{Complex Representation}

\textbf{Experimental density data.}
We curate X-ray protein--ligand instances from PLINDER
\citep{durairaj2024plinder} and retrieve the experimental $2mF_o-DF_c$ map
associated with each parent PDB entry. We retain valid, noncovalently bound
ligands that are fully resolved, contain 6--50 heavy atoms, and have a ligand
real-space correlation coefficient (RSCC) of at least $0.8$. We exclude all PDB
entries used for downstream validation or testing, yielding more than 100K
pre-training instances. The complete curation and normalization procedures are
given in \cref{appx:data-detail}.

\textbf{Complex voxel representation.}
We construct a binding-site grid centered at the ligand centroid, following
\citet{pinheiro2024structure}. The grid contains $G^3$ voxels with $G=64$ and
$0.25$\,\AA{} spacing, covering a $16$\,\AA{} cube. Each complex is represented
as follows:
\begin{equation}
\mathbf{x}=[\mathbf{a};\rho;\lVert\nabla\rho\rVert]
\in\mathbb{R}^{G^3\times13},
\qquad
a_{v,c}=1-\prod_{i\in\mathcal{A}_c}(1-g_{v,i}),
\qquad
g_{v,i}=\exp\!\left(
-\frac{\lVert\mathbf r_v-\mathbf r_i\rVert_2^2}{(0.93R_i)^2}
\right).
\label{eq:complex-voxel}
\end{equation}
Here, $\mathbf a$ contains seven ligand (C, O, N, S, F, Cl, P) and four pocket
(C, O, N, S) atom-type channels. The set $\mathcal A_c$ contains atoms assigned
to channel $c$, and $a_{v,c}$ combines their Gaussian occupancies at voxel $v$.
The vectors $\mathbf r_v$ and $\mathbf r_i$ denote the voxel-center and atom
positions, respectively, and $R_i$ is the element-specific van der Waals radius
of atom $i$. We resample the crystallographic map onto the same grid by
trilinear interpolation and compute its gradient magnitude on the resampled
grid \citep{suriana2025beyond}. Further voxelization and normalization details
are provided in \cref{appx:method-detail}.

\subsection{Density-Aware Masked Pre-training}

We divide the grid into non-overlapping blocks of $P^3$ voxels, with $P=8$, and
embed the ligand, pocket, and density channel groups using a Channel-ViT
\citep{bao2023channel-24a}. Following masked autoencoding
\citep{he2022masked}, we mask $75\%$ of the spatial blocks, biasing selection
toward atom-containing regions, and apply the same spatial mask to every
channel. The encoder $f_\theta$ and a lightweight reconstruction head $h_\psi$
predict all channels at the masked voxels:
\begin{equation}
\widehat{\mathbf x}=h_\psi\!\left(f_\theta(\widetilde{\mathbf x})\right),
\qquad
\mathcal L(\theta,\psi)=
\frac{1}{|\mathcal M|C}
\sum_{v\in\mathcal M}\sum_{c=1}^{C}
w_c\,\omega_{v,c}
\left(\widehat{x}_{v,c}-x_{v,c}\right)^2.
\label{eq:masked-pretraining}
\end{equation}
Here, $\widetilde{\mathbf x}$ is the masked input, $\mathcal M$ denotes the
masked voxels, and $C=13$ is the number of channels. The channel weight $w_c$
balances different occupancy frequencies, while $\omega_{v,c}$ gives additional
weight to occupied voxels in sparse atomic channels. Architecture, weighting,
augmentation, and optimization details are provided in
\cref{appx:method-detail}.

\subsection{Downstream Transfer}

\textbf{Binding-affinity regression.}
We pass the unmasked 13-channel complex grid through the frozen encoder and
mean-pool its final tokens. A two-layer MLP maps the pooled representation to
binding affinity. Only the probe is trained, so the task directly evaluates the
transferred representation.

\textbf{Reference-assisted molecular generation.}
We augment VoxBind with the frozen encoder's spatial features
\citep{pinheiro2024structure}. \ours{} processes the reference ligand, protein
pocket, and experimental map, producing one feature for each spatial block. We
broadcast each feature to its component voxels and fuse it with VoxBind through
a trainable, zero-initialized residual projection \citep{zhang2023adding}.
\ours{} remains frozen, while the projection and generative model are trained;
the remaining denoising architecture and sampling procedure follow VoxBind.
